% s-037-p037-m002; original PDF 37, 38, 39
\Aphorism{s-037-p037-m002}{23}
Das „Recht zur Hekatombe“, jenes urälteste Philosophenvorrecht, niedrige Arten von Glück und Leben zu Gunsten der höheren zu opfern, setzt zu allererst ein \Sperr{gutes Gewissen} voraus, und zwar eines, das man wirklich hat und sich nicht durch theoretischen Radicalismus und allerhand künstlichen Lärm einzureden braucht.
\par

In dieser Lage sind z. B. wir Modernen; wenn wir die Probleme einer socialen Rangordnung erwägen, so denken wir „wider den Strich“ und müssen durch wilde Paradoxenmacherei über den dumpfen Widerstand unserer Instincte Herr werden: aus unserer allzuberedsamen Vergötterung des souveränen Ich ist deutlich genug das noch unüberwundene schlechte Gewissen herauszuhören. Wir fleisch- und blutgewordenen Ideale von 1789, die wir mit Herz, Nieren und Rückenmark an das gleiche Recht Aller glauben und damit jeden cerebralen Gegenglauben niederstimmen, wir bedürfen ungeheurer Selbstverleugnung und Moralität, um uns das Zugeständniss vom „Vorrecht der Wenigsten“ abzuzwingen — als Zugeständniss, noch nicht einmal als selbstgeerntete Erfahrung, und noch lange nicht als Glauben, als Axiom, als vererbte Gewissheit vor aller Erfahrung. Und doch finge erst damit jenes Vorrecht in Wirklichkeit an, und wo solch ein Glaube fehlt, ist es zum aristokratischen Individualismus noch um ein paar Generationen zu früh. Das ist der Fall Raskolnikow, mit dessen Aufdeckung Dostojewskij die Ohnmacht jeder blossen Theorie, jeder vernünftigen oder wissenschaftlichen Ueberzeugung blosslegt, für die nicht der \Sperr{Leib} Partei ergreift: das Symptom einer solchen Einverleibung, die Garantie des guten Einvernehmens zwischen Wille und Intellect ist eben der Gefühlston Lust, das Freiwerden der Bahn des geringsten Widerstandes, die Association mit dem guten Gewissen. Die theoretische Gewöhnung an bestimmte Denkweisen ist nur Vorstufe und Mittelglied dieser Entwicklung, eine Väter- und Grossväterarbeit zu Gunsten der Enkel: Raskolnikow hätte Sonja heirathen und seinen Sohn in einer suggestiven Gottesgnaden-Atmosphäre grossziehen müssen — dieser oder einer seiner Nachkommen wäre dann vielleicht der Napoleon geworden, der ohne mit der Wimper zu zucken Tausende aufs Schlachtfeld führt. Es ist eine tiefe Weisheit, dass man ohne Ahnen nicht legitim sein kann; wer die Zukunft vorweglebt, dessen wird leicht die Moral, schlimmer noch, die Justiz der Gegenwart Herr.
\par

